For axisymmetric plasmas, The Drift-Kinetic Equation is given by 
\begin{equation}
\frac{\partial f}{\partial t}+{\bf v}_{\text{gc}}\cdot \nabla f={\cal C}%
\left( f\right) +{\cal Q}\left( f\right) +{\cal E}\left( f\right) 
\label{dke1}
\end{equation}
where $f=f\left( p,\xi ,\psi ,\theta ,t\right) $ is the guiding-center
distribution function, and the guiding center velocity is decomposed into
the fast parellel motion along the field lines, and a drift velocity ${\bf v}%
_{D}$. 
\begin{equation}
{\bf v}_{\text{gc}}=v_{\Vert }\widehat{b}+{\bf v}_{D}  \label{vgc}
\end{equation}

The magnetic field is generally expressed as 
\begin{equation}
{\bf B}=I\left( \psi \right) \nabla \phi +\nabla \psi \times \nabla \phi
\label{bexp}
\end{equation}
which becomes; in $\left( \psi ,s,\phi \right) $ coordinates, 
\begin{equation}
{\bf B}=\frac{I\left( \psi \right) }{R}\widehat{\phi }-\frac{\left| \nabla
\psi \right| }{R}\widehat{s}  \label{bexp2}
\end{equation}

The gradient in $\left( \psi ,s,\phi \right) $ coordinates is 
\begin{eqnarray}
\nabla &=&\nabla \psi \frac{\partial }{\partial \psi }+\nabla s\frac{%
\partial }{\partial s}+\nabla \phi \frac{\partial }{\partial \phi } \\
&=&\nabla \psi \frac{\partial }{\partial \psi }+\widehat{s}\frac{\partial }{%
\partial s}+\frac{\widehat{\phi }}{R}\frac{\partial }{\partial \phi }
\label{gradient}
\end{eqnarray}

The constants of the motion are the total energy 
\begin{equation}
\varepsilon =\frac{1}{2}m\left( v_{\Vert }^{2}+v_{\bot }^{2}\right) 
\end{equation}
and the magnetic moment 
\begin{equation}
\mu =\frac{mv_{\bot }^{2}}{2B}  \label{mudef}
\end{equation}
so that we have the conservations laws 
\begin{eqnarray}
\frac{\partial \mu }{\partial s} &=&0  \label{mucons} \\
\frac{\partial }{\partial s}\left[ v_{\Vert }^{2}+\frac{2\mu B}{m}\right] 
&=&0  \label{encons}
\end{eqnarray}

\subsection{Drift Velocity from the Conservation of Canonical Momentum}

The toroidal canonical momentum is also a constant of the motion because of
axisymmetry. It is expressed as 
\begin{equation}
P_{\phi }=R\left[ \gamma mv_{\phi }+q_{e}A_{\phi }\right]
\end{equation}
where $A_{\phi }$ is the toroidal component of the vector potential. From
the relation 
\begin{equation}
{\bf B}=\nabla \times {\bf A}
\end{equation}
and the expression (\ref{curlpsis}) of a rotational in $\left( \psi ,s,\phi
\right) $ coordinates, we get 
\begin{eqnarray}
{\bf B} &=&\left[ \frac{1}{R}\frac{\partial }{\partial s}\left( RA_{\phi
}\right) -\frac{1}{R}\frac{\partial }{\partial \phi }\left( A_{s}\right)
\right] \widehat{\psi } \\
&&+\left[ \frac{1}{R}\frac{\partial }{\partial \phi }\left( A_{\psi }\right)
-\frac{\left| \nabla \psi \right| }{R}\frac{\partial }{\partial \psi }\left(
RA_{\phi }\right) \right] \widehat{s} \\
&&+\left[ \left| \nabla \psi \right| \frac{\partial }{\partial \psi }\left(
A_{s}\right) -\left| \nabla \psi \right| \frac{\partial }{\partial s}\left( 
\frac{A_{\psi }}{\left| \nabla \psi \right| }\right) \right] \widehat{\phi }
\end{eqnarray}
with 
\begin{eqnarray}
A_{\psi } &=&{\bf A}\cdot \widehat{\psi } \\
A_{s} &=&{\bf A}\cdot \widehat{s} \\
A_{\phi } &=&{\bf A}\cdot \widehat{\phi }
\end{eqnarray}

In axisymmetric plasma, this reduces to 
\begin{eqnarray}
{\bf B} &=&\frac{1}{R}\frac{\partial }{\partial s}\left( RA_{\phi }\right) 
\widehat{\psi }  \label{bpsi} \\
&&-\frac{\left| \nabla \psi \right| }{R}\frac{\partial }{\partial \psi }%
\left( RA_{\phi }\right) \widehat{s} \\
&&+\left[ \left| \nabla \psi \right| \frac{\partial }{\partial \psi }\left(
A_{s}\right) -\left| \nabla \psi \right| \frac{\partial }{\partial s}\left( 
\frac{A_{\psi }}{\left| \nabla \psi \right| }\right) \right] \widehat{\phi }
\end{eqnarray}
so that 
\begin{equation}
B_{s}=-\frac{\left| \nabla \psi \right| }{R}\frac{\partial }{\partial \psi }%
\left( RA_{\phi }\right)
\end{equation}

In addition, we know from (\ref{bexp2}) that 
\begin{equation}
B_{s}=-\frac{\left| \nabla \psi \right| }{R}
\end{equation}
so that be obtain 
\begin{equation}
\frac{\partial RA_{\phi }}{\partial \psi }=1  \label{raphi}
\end{equation}

Because the toroidal canonical momentum is a constant of the motion, we have
that \thinspace 
\begin{equation}
{\bf v}_{\text{gc}}\cdot \nabla P_{\phi }=0
\end{equation}
which can be decomposed into 
\begin{equation}
{\bf v}_{\text{gc}}\cdot \nabla \left( R\gamma mv_{\phi }\right) +{\bf v}_{%
\text{gc}}\cdot \nabla \left( q_{e}RA_{\phi }\right) =0  \label{eqpphi}
\end{equation}

Using (\ref{psisgrad}), we get 
\begin{equation}
{\bf v}_{\text{gc}}\cdot \nabla \left( qRA_{\phi }\right) ={\bf v}_{\text{gc}%
}\cdot \left[ \nabla \psi \frac{\partial }{\partial \psi }+\widehat{s}\frac{%
\partial }{\partial s}+\frac{\widehat{\phi }}{R}\frac{\partial }{\partial
\phi }\right] \left( qRA_{\phi }\right)
\end{equation}
which in axisymmetric systems gives 
\begin{equation}
{\bf v}_{\text{gc}}\cdot \nabla \left( qA_{\phi }\right) =q{\bf v}_{\text{gc}%
}\cdot \left[ \nabla \psi \frac{\partial RA_{\phi }}{\partial \psi }+%
\widehat{s}\frac{\partial RA_{\phi }}{\partial s}\right]
\end{equation}

Since $B_{\psi }=0$, we have from (\ref{bpsi}) that $\partial \left(
RA_{\phi }\right) /\partial s=0$ and therefore, using (\ref{raphi}), 
\begin{equation}
{\bf v}_{\text{gc}}\cdot \nabla \left( qA_{\phi }\right) =q{\bf v}_{\text{gc}%
}\cdot \nabla \psi
\end{equation}

The only velocity accross the flux-surfaces is the drift we are looking for,
so that we get, using (\ref{vgc}) 
\begin{equation}
{\bf v}_{\text{gc}}\cdot \nabla \left( qA_{\phi }\right) =q{\bf v}_{D}\cdot
\nabla \psi
\end{equation}
and the equation (\ref{eqpphi}) becomes 
\begin{equation}
q{\bf v}_{D}\cdot \nabla \psi =-{\bf v}_{\text{gc}}\cdot \nabla \left(
R\gamma mv_{\phi }\right)
\end{equation}

Assuming a priori that $\left| v_{\Vert }\right| \gg \left| {\bf v}%
_{D}\right| $, this equation reduces to 
\begin{eqnarray}
{\bf v}_{D}\cdot \nabla \psi &=&-\frac{1}{q}\frac{v_{\Vert }}{B}{\bf B}\cdot
\nabla \left( R\gamma mv_{\phi }\right) \\
&=&-\frac{v_{\Vert }}{\Omega }{\bf B}\cdot \nabla \left( Rv_{\phi }\right)
\end{eqnarray}
where we used that $\partial \gamma /\partial s=0$ because of the
conservation of energy

The toroidal velocity is related to the parallel velocity by 
\begin{equation}
v_{\phi }=\frac{B_{\phi }}{B}v_{\Vert }=\frac{I\left( \psi \right) }{RB}%
v_{\Vert }
\end{equation}
so that 
\begin{equation}
Rv_{\phi }=\frac{I\left( \psi \right) }{B}v_{\Vert }
\end{equation}

Since $I\left( \psi \right) $ is a flux function, it can be taken out of the
gradient, so that 
\begin{equation}
{\bf v}_{D}\cdot \nabla \psi =-\frac{v_{\Vert }}{\Omega }I\left( \psi
\right) {\bf B}\cdot \nabla \left( \frac{v_{\Vert }}{B}\right)
\label{vdpsi1}
\end{equation}

\subsection{Drift Velocity from the Expression of Single Particle Drift}

The guiding-center drift velocity due to the magnetic field gradient and
curvature is

\begin{equation}
{\bf v}_{D}=\frac{1}{\Omega }\left( v_{\Vert }^{2}+\frac{v_{\bot }^{2}}{2}%
\right) \frac{{\bf B}\times \nabla B}{B^{2}}
\end{equation}
and its component perpendicular to the flux-surface can be written as 
\begin{equation}
{\bf v}_{D}\cdot \nabla \psi =\frac{1}{\Omega }\left( v_{\Vert }^{2}+\frac{%
v_{\bot }^{2}}{2}\right) \frac{1}{B^{2}}\nabla \psi \times {\bf B}\cdot
\nabla B
\end{equation}

Inserting the expression (\ref{bexp2}) of the magnetic field, we find 
\begin{equation}
{\bf v}_{D}\cdot \nabla \psi =-\frac{1}{\Omega }\left( v_{\Vert }^{2}+\frac{%
v_{\bot }^{2}}{2}\right) \frac{\left| \nabla \psi \right| }{B^{2}R}\left[
I\left( \psi \right) \widehat{s}+\left| \nabla \psi \right| \widehat{\phi }%
\right] \cdot \nabla B  \label{drift}
\end{equation}

Using (\ref{gradient}), the equation (\ref{drift}) becomes 
\begin{equation}
{\bf v}_{D}\cdot \nabla \psi =-\frac{1}{\Omega }\left( v_{\Vert }^{2}+\frac{%
v_{\bot }^{2}}{2}\right) \frac{\left| \nabla \psi \right| }{B^{2}R}\left[
I\left( \psi \right) \frac{\partial B}{\partial s}+\frac{\left| \nabla \psi
\right| }{R}\frac{\partial B}{\partial \phi }\right]
\end{equation}

Under the assumption of axisymmetry, we are left with 
\begin{equation}
{\bf v}_{D}\cdot \nabla \psi =-\frac{1}{\Omega }\left( v_{\Vert }^{2}+\frac{%
v_{\bot }^{2}}{2}\right) \frac{\left| \nabla \psi \right| I\left( \psi
\right) }{B^{2}R}\frac{\partial B}{\partial s}
\end{equation}

With the definition (\ref{mudef }) of the magnetic moment $\mu $, we rewrite 
\begin{equation}
{\bf v}_{D}\cdot \nabla \psi =-\frac{1}{\Omega }\frac{\left| \nabla \psi
\right| I\left( \psi \right) }{B^{2}R}\left( v_{\Vert }^{2}+\frac{\mu B}{m}%
\right) \frac{\partial B}{\partial s}  \label{drift2}
\end{equation}
We have, using the conservation of magnetic momentum (\ref{mucons}), 
\begin{equation}
\left( v_{\Vert }^{2}+\frac{\mu B}{m}\right) \frac{\partial B}{\partial s}%
=v_{\Vert }^{2}\frac{\partial B}{\partial s}+B\frac{\partial }{\partial s}
\left( \frac{\mu B}{m}\right)
\end{equation}

Using the conservation of energy (\ref{encons}), we get 
\begin{eqnarray}
\left( v_{\Vert }^{2}+\frac{\mu B}{m}\right) \frac{\partial B}{\partial s}
&=&v_{\Vert }^{2}\frac{\partial B}{\partial s}-B\frac{\partial }{\partial s}%
\left( \frac{v_{\Vert }^{2}}{2}\right) \\
&=&-v_{\Vert }\left[ B\frac{\partial v_{\Vert }}{\partial s}-v_{\Vert }\frac{%
\partial B}{\partial s}\right] \\
&=&-v_{\Vert }B^{2}\frac{\partial }{\partial s}\left( \frac{v_{\Vert }}{B}%
\right)
\end{eqnarray}
and finally, the equation (\ref{drift2}) becomes 
\begin{equation}
{\bf v}_{D}\cdot \nabla \psi =\frac{v_{\Vert }}{\Omega }I\left( \psi \right) 
\frac{\left| \nabla \psi \right| }{R}\frac{\partial }{\partial s}\left( 
\frac{v_{\Vert }}{B}\right)  \label{drift3}
\end{equation}

In addition, 
\begin{eqnarray}
{\bf B}\cdot \nabla &=&I\left( \psi \right) \nabla \phi \cdot \nabla +\nabla
\psi \times \nabla \phi \cdot \nabla \\
&=&\frac{I\left( \psi \right) }{R}\frac{\partial }{\partial \phi }-\frac{%
\left| \nabla \psi \right| }{R}\frac{\partial }{\partial s}
\end{eqnarray}
and, using axisymmetry, 
\begin{equation}
{\bf B}\cdot \nabla =-\frac{\left| \nabla \psi \right| }{R}\frac{\partial }{%
\partial s}
\end{equation}
so that we can rewrite (\ref{drift3}) as 
\begin{equation}
{\bf v}_{D}\cdot \nabla \psi =-\frac{v_{\Vert }}{\Omega }I\left( \psi
\right) {\bf B}\cdot \nabla \left( \frac{v_{\Vert }}{B}\right)
\label{vdpsi2}
\end{equation}

\subsection{Case of Circular concentric flux-surfaces}

In this case, $\psi =$ $\psi \left( r\right) $ and therefore 
\begin{equation}
\nabla \psi =\psi ^{\prime }\left( r\right) \widehat{r}
\end{equation}
and 
\begin{equation}
v_{Dr}={\bf v}_{D}\cdot \widehat{r}=\frac{{\bf v}_{D}\cdot \nabla \psi }{%
\psi ^{\prime }}
\end{equation}
which gives 
\begin{equation}
v_{Dr}=-\frac{v_{\Vert }}{\Omega }\frac{I\left( \psi \right) }{\psi ^{\prime
}}{\bf B}\cdot \nabla \left( \frac{v_{\Vert }}{B}\right)
\end{equation}

In addition, 
\begin{equation}
{\bf B}\cdot \nabla =-\frac{\left| \nabla \psi \right| }{R}\frac{\partial }{
\partial s}=-\frac{\left| \psi ^{\prime }\right| }{R}\frac{\partial }{%
\partial s}
\end{equation}
and, because 
\begin{equation}
\widehat{s}\cdot \widehat{\theta }=\widehat{\psi }\cdot \widehat{r}=\sigma
_{\psi }
\end{equation}
we find 
\begin{equation}
\frac{\partial }{\partial s}=\frac{\sigma _{\psi }}{r}\frac{\partial }{%
\partial \theta }
\end{equation}
and 
\begin{equation}
{\bf B}\cdot \nabla =-\frac{\psi ^{\prime }}{Rr}\frac{\partial }{\partial
\theta }
\end{equation}
so that finally 
\begin{equation}
v_{Dr}=\frac{v_{\Vert }}{\Omega }\frac{I\left( \psi \right) }{Rr}\frac{
\partial }{\partial \theta }\left( \frac{v_{\Vert }}{B}\right) =\frac{%
v_{\Vert }}{r}\frac{I\left( \psi \right) }{RB}\frac{\partial }{\partial
\theta }\left( \frac{v_{\Vert }}{\Omega }\right)
\end{equation}

When the toroidal field dominates, $B\simeq I\left( \psi \right) /R$ and 
\begin{equation}
v_{Dr}\simeq \frac{v_{\Vert }}{r}\frac{\partial }{\partial \theta }\left( 
\frac{v_{\Vert }}{\Omega }\right)
\end{equation}

\subsection{Steady-State Drift-Kinetic Equation}

Using (\ref{dke1}), (\ref{vgc}) and (\ref{vdpsi1}) or (\ref{vdpsi2}), we
obtain in steady-state 
\begin{equation}
v_{\Vert }\widehat{b}\cdot \nabla f-\frac{v_{\Vert }}{\Omega }I\left( \psi
\right) {\bf B}\cdot \nabla \left( \frac{v_{\Vert }}{B}\right) \frac{%
\partial f}{\partial \psi }={\cal C}\left( f\right) +{\cal Q}\left( f\right)
+{\cal E}\left( f\right) 
\end{equation}
which can be rewritten as 
\begin{equation}
v_{s}\frac{\partial f}{\partial s}+\frac{v_{\Vert }}{\Omega }I\left( \psi
\right) \frac{\left| \nabla \psi \right| }{R}\frac{\partial }{\partial s}%
\left( \frac{v_{\Vert }}{B}\right) \frac{\partial f}{\partial \psi }={\cal C}%
\left( f\right) +{\cal Q}\left( f\right) +{\cal E}\left( f\right) 
\end{equation}
with 
\begin{equation}
v_{s}=v_{\Vert }\left( \widehat{b}\cdot \widehat{s}\right) 
\end{equation}
